\section{Konstruksi Recursive Descent yang Sinkron dengan Tabel}
\label{sec:ch10-implementasi}

\subsection{Dari Tabel ke Fungsi}

Setiap nonterminal menjadi satu fungsi. Isi fungsi mengikuti aturan
pada tabel LL(1):
\begin{itemize}[leftmargin=*]
  \item untuk aturan \(A \to X_1\cdots X_k\), panggil
        \code{parseX1}, \code{parseX2}, dan seterusnya;
  \item untuk terminal, panggil \code{match};
  \item untuk \(\varepsilon\), tidak ada aksi;
  \item untuk sel kosong, laporkan kesalahan.
\end{itemize}

Tabel berikut memetakan nonterminal ke fungsi parser:
\begin{table}[H]
\centering
\caption{Pemetaan nonterminal ke fungsi parser}
\label{tab:ch10-fungsi}
\begin{tabularx}{\textwidth}{@{}l l X@{}}
\toprule
\textbf{Nonterminal} & \textbf{Fungsi} & \textbf{Tanggung jawab} \\
\midrule
\code{program} & \code{parser\_program} & seluruh berkas \\
\code{topdecl} & \code{parse\_topdecl} & deklarasi fungsi \\
\code{block} & \code{parse\_block} & daftar pernyataan \\
\code{stmt} & \code{parse\_stmt} & satu pernyataan \\
\code{expr}--\code{primary} & \code{parse\_expr}, ... &
ekspresi berprioritas \\
\bottomrule
\end{tabularx}
\end{table}

\subsection{Pemilihan Berdasarkan Token}

Fungsi \code{parse\_stmt} memilih cabang berdasarkan token awal.
Potongan implementasinya ditunjukkan pada
\lstref{lst:ch10-stmt}.

\includecode[
  language=C,
  firstline=296,
  lastline=340,
  firstnumber=286,
  caption={Bagian awal \code{parse\_stmt}},
  label={lst:ch10-stmt}
]{../../code/subsetc/src/parser.c}

Setiap cabang memeriksa jenis token dan memanggil fungsi yang sesuai.
Jika tidak ada cabang yang cocok, parser melaporkan kesalahan dan
mengembalikan simpul pengganti.

\subsection{Sinkronisasi}

Setelah kesalahan, parser harus kembali ke jalur yang dapat
dilanjutkan. Fungsi \code{synchronize} memajukan token sampai titik
aman, misalnya titik koma atau kurung tutup. Implementasinya
ditunjukkan pada \lstref{lst:ch10-sync}.

\includecode[
  language=C,
  firstline=42,
  lastline=55,
  firstnumber=42,
  caption={Fungsi \code{synchronize}},
  label={lst:ch10-sync}
]{../../code/subsetc/src/parser.c}

Titik sinkronisasi yang umum adalah \code{;}, \code{\}}, dan awal
pernyataan berikutnya. Pemilihan titik sinkronisasi memengaruhi
jumlah kesalahan yang dilaporkan: terlalu agresif dapat menghasilkan
laporan palsu, terlalu konservatif dapat melewatkan kesalahan.

\subsection{Verifikasi dengan Tabel}

Setelah implementasi, bandingkan perilaku parser dengan tabel LL(1):
\begin{enumerate}[label=\arabic*.,leftmargin=*]
  \item untuk setiap sel tabel, uji masukan yang memicu aturan
        tersebut;
  \item pastikan parser tidak memilih aturan yang salah;
  \item uji sel kosong dan pastikan kesalahan dilaporkan.
\end{enumerate}

Verifikasi ini penting karena recursive descent yang ditulis manual
dapat menyimpang dari tabel tanpa terlihat. Pengujian sistematis
membantu menemukan perbedaan tersebut.
